\section{Постановка задачи}
\label{sec:problem}

В настоящем разделе формулируется содержательная и формальная постановка
задачи, определяются объект и предмет исследования, описываются требования
к разрабатываемой системе и ограничения, накладываемые на решение.

\subsection{Содержательная постановка задачи}

Требуется разработать программное обеспечение мобильного робота на
омни-платформе с колёсами типа Mecanum, оснащённого фронтальной
RGB-камерой и двумерным лазерным дальномером (2D-лидаром), которое в
условиях типовой квартирной среды с подвижными препятствиями обеспечивает:
\begin{enumerate}
    \item автономное обнаружение подвижного пользователя в поле зрения
          камеры;
    \item сопровождение пользователя с удержанием целевой дистанции до
          него в заданных пределах;
    \item удержание пользователя в центральной зоне поля зрения камеры
          по горизонтальной оси;
    \item избегание столкновений со статическими и подвижными
          препятствиями;
    \item безопасное поведение при временной потере наблюдаемости цели
          (остановка с последующим поиском).
\end{enumerate}

Среда эксплуатации --- квартира с характерными габаритами и типичным
набором бытовых препятствий. Среда считается частично известной: карта
помещения формируется в реальном времени средствами SLAM по данным лидара.

\subsection{Объект и предмет исследования}

\textbf{Объект исследования:} автономный мобильный колёсный робот на
омни-платформе, предназначенный для задачи сопровождения пользователя в
ограниченной помещении.

\textbf{Предмет исследования:} алгоритмы и программные средства
локального планирования траектории и управления движением такого робота
в условиях динамически меняющейся среды.

\subsection{Формальная постановка задачи управления}

Вектор состояния робота в глобальной системе
координат определяется как
\begin{equation}
\mathbf{x}_k = \begin{bmatrix} x_k & y_k & \psi_k \end{bmatrix}^{\top},
\label{eq:state}
\end{equation}
где $x_k, y_k$~--- координаты центра масс робота в глобальной системе,
$\psi_k$~--- угол ориентации корпуса относительно оси $x$ глобальной
системы координат.

Вектор управления, доступный для омни-платформы, имеет вид
\begin{equation}
\mathbf{u}_k = \begin{bmatrix} v_{x,k} & v_{y,k} & \omega_k \end{bmatrix}^{\top},
\label{eq:control}
\end{equation}
где $v_{x,k}, v_{y,k}$~--- составляющие линейной скорости корпуса в связанной
с роботом системе координат $x_B y_B$, $\omega_k$~--- угловая скорость
вращения корпуса относительно вертикальной оси.

Положение подвижного пользователя (цели) в глобальной системе координат
описывается вектором
\begin{equation}
\mathbf{x}^{t}_k = \begin{bmatrix} x^{t}_k & y^{t}_k \end{bmatrix}^{\top}.
\label{eq:target_state}
\end{equation}

Пусть $d_k = \|\mathbf{p}_k - \mathbf{x}^{t}_k\|_2$~--- расстояние от робота до
цели, где $\mathbf{p}_k = [x_k, y_k]^{\top}$. Пусть $\phi_k$~--- угол
направления на цель относительно оси камеры. Тогда задача управления
формулируется как поиск последовательности управлений $\{\mathbf{u}_k\}$,
минимизирующей составную функцию потерь
\begin{equation}
J = \sum_{k=0}^{N-1} \ell(\mathbf{x}_k, \mathbf{u}_k, \mathbf{x}^{t}_k)
  + \ell_f(\mathbf{x}_N),
\label{eq:cost_general}
\end{equation}
при ограничениях модели движения $\mathbf{x}_{k+1} = f(\mathbf{x}_k,\mathbf{u}_k)$,
ограничениях на управление $\mathbf{u}_k \in \mathcal{U}$ и условиях
безопасности
\begin{equation}
\|\mathbf{p}_k - \mathbf{p}^{\text{obs}}_{i,k}\|_2 \ge d_{\min}, \quad \forall i, \forall k,
\label{eq:safety}
\end{equation}
где $\mathbf{p}^{\text{obs}}_{i,k}$~--- положение $i$-го препятствия на шаге $k$,
$d_{\min}$~--- минимально допустимое расстояние.

Конкретный вид функции $\ell(\cdot)$, отражающий прикладные требования, и
детализация модели $f(\cdot)$ приведены в главе~\ref{sec:solution}.

\subsection{Требования к системе}

Разрабатываемая система должна удовлетворять следующим требованиям:

\subsubsection{Функциональные требования}

\begin{enumerate}
    \item Автономное обнаружение цели в кадре камеры по характерному
          визуальному признаку (в рамках модельной реализации~--- по
          цветовой сегментации).
    \item Пересчёт положения обнаруженной цели в карту помещения с
          использованием данных лидара, дерева преобразований координат и координат цели в кадре камеры.
    \item Удержание целевой дистанции $d^{\star}$ до пользователя в
          пределах заданной погрешности.
    \item Удержание цели в центральной зоне кадра камеры по
          горизонтальной оси.
    \item Обход статических и динамических препятствий с соблюдением
          минимальной дистанции $d_{\min}$.
    \item Поиск цели в случае ее потери.
\end{enumerate}

\subsubsection{Нефункциональные требования}

\begin{enumerate}
    \item Воспроизводимость среды запуска за счёт контейнеризации.
    \item Возможность запуска всей симуляционной цепочки одной командой.
    \item Модульность архитектуры: перенос алгоритмических модулей на
          реальное оборудование без изменения их внутренней логики.
    \item Параметризуемость поведения системы через конфигурационные
          файлы без необходимости перекомпиляции.
    \item Удалённый доступ к графическому интерфейсу симулятора
          независимо от особенностей графической подсистемы хоста.
\end{enumerate}

\subsubsection{Ограничения}

\begin{enumerate}
    \item Работа выполняется в рамках физически-реалистичной симуляции
          Gazebo, интеграция на реальное оборудование выходит за рамки
          данной ВКР.
    \item Пользователь представлен моделью подвижного цветного объекта
          (синий цилиндр), что позволяет сконцентрироваться на контуре
          управления; замена детектора цели на промышленное решение
          компьютерного зрения технически осуществима, но не включена в
          объём работы.
    \item Сенсорный набор ограничен 2D-лидаром и фронтальной
          RGB-камерой.
\end{enumerate}

\subsection{Критерии оценки качества решения}

Качество разработанного решения оценивается следующими
количественными характеристиками, вычисляемыми по записанным
rosbag2 (подробное описание методики и формальные определения~---
в главе~\ref{sec:experiments}):
\begin{itemize}
    \item доля времени каждого режима сопровождения~$\bar{\tau}_m$~---
          доминирование режима \textit{Сопровождение цели}
          подтверждает работоспособность связки восприятия и
          контроллера;
    \item доля времени видимости цели $\bar{\nu} \in [0,1]$~---
          интегральная характеристика устойчивости трекера и
          поведения при загораживаниях;
    \item число фактических контактов с препятствиями
          $N_{\text{c}}$ и число заходов в зону риска
          $N_{\text{r}}$, а также доля предотвращённых контактов
          $\rho = N_{\text{p}}/N_{\text{r}}$~--- характеристики
          безопасности относительно ограничения~\eqref{eq:safety};
    \item среднеквадратичная норма приращения управления
          $\mathrm{RMS}(\|\Delta\mathbf{u}_k\|_2)$~--- характеристика
          плавности управления.
\end{itemize}
